\chapter{Документация}
\label{ch:docs}

\hypertarget{ux432ux432ux435ux434ux435ux43dux438ux435}{%
\section{Введение}\label{ux432ux432ux435ux434ux435ux43dux438ux435}}

\hypertarget{ux43eux43fux438ux441ux430ux43dux438ux435-ux43fux440ux43eux435ux43aux442ux430}{%
\subsection*{Описание
проекта}\label{ux43eux43fux438ux441ux430ux43dux438ux435-ux43fux440ux43eux435ux43aux442ux430}}

Проект CinemaPon представляет собой бэкенд-систему,
предназначенную для управления данными о фильмах и связанными с ними
сущностями, такими как жанры, режиссеры, теги, рейтинги и комментарии.
Основная цель системы --- предоставить удобный инструмент для хранения,
обработки и анализа информации о фильмах, а также для генерации
персонализированных рекомендаций для пользователей.

Система поддерживает стандартные CRUD-операции (Create, Read, Update,
Delete), что позволяет администраторам легко управлять данными. Кроме
того, в проект интегрирован модуль машинного обучения (ML), который
ежедневно обновляет рекомендации для пользователей и фильмов на основе
их предпочтений и поведения.

\hypertarget{ux446ux435ux43bux435ux432ux430ux44f-ux430ux443ux434ux438ux442ux43eux440ux438ux44f}{%
\subsection*{Целевая
аудитория}\label{ux446ux435ux43bux435ux432ux430ux44f-ux430ux443ux434ux438ux442ux43eux440ux438ux44f}}

Целевая аудитория проекта делится на две основные группы: 
\begin{enumerate}
    \item 
Администраторы системы --- лица, ответственные за управление
данными о фильмах, жанрах, режиссерах, тегах и других сущностях. Они
имеют полный доступ ко всем функциям системы, включая добавление,
редактирование и удаление данных. 
    \item 
Пользователи --- конечные
потребители, которые взаимодействуют с системой через API для получения
информации о фильмах, рекомендаций, а также для оставления комментариев
и оценок.
\end{enumerate}
\hypertarget{ux43eux441ux43dux43eux432ux43dux44bux435-ux444ux443ux43dux43aux446ux438ux438}{%
\subsection*{Основные
функции}\label{ux43eux441ux43dux43eux432ux43dux44bux435-ux444ux443ux43dux43aux446ux438ux438}}

Система предоставляет следующие ключевые функции: 
\begin{itemize}
    \item 
Управление
данными о фильмах: Добавление, редактирование, удаление и просмотр
информации о фильмах.
\item
Управление связанными сущностями:
Работа с жанрами, режиссерами, тегами, рейтингами и комментариями. 
\item
Генерация рекомендаций: Ежедневное обновление рекомендаций для
пользователей и фильмов с использованием машинного обучения. 
\item
Аутентификация и авторизация: Поддержка регистрации, входа в
систему и управления правами доступа.
\end{itemize}

\hypertarget{ux430ux440ux445ux438ux442ux435ux43aux442ux443ux440ux430-ux441ux438ux441ux442ux435ux43cux44b}{%
\section{Архитектура
системы}\label{ux430ux440ux445ux438ux442ux435ux43aux442ux443ux440ux430-ux441ux438ux441ux442ux435ux43cux44b}}

\hypertarget{ux43eux431ux449ux430ux44f-ux430ux440ux445ux438ux442ux435ux43aux442ux443ux440ux430}{%
\subsection*{Общая
архитектура}\label{ux43eux431ux449ux430ux44f-ux430ux440ux445ux438ux442ux435ux43aux442ux443ux440ux430}}

Система построена на основе микросервисной архитектуры, где бэкенд
реализован с использованием Django и Django REST
Framework (DRF). Это позволяет системе быть масштабируемой и легко
расширяемой. Модуль машинного обучения интегрирован в систему и работает
в фоновом режиме, обновляя рекомендации ежедневно.

\hypertarget{ux442ux435ux445ux43dux43eux43bux43eux433ux438ux447ux435ux441ux43aux438ux439-ux441ux442ux435ux43a}{%
\subsection*{Технологический
стек}\label{ux442ux435ux445ux43dux43eux43bux43eux433ux438ux447ux435ux441ux43aux438ux439-ux441ux442ux435ux43a}}

\begin{itemize}
\tightlist
\item
  Бэкенд: Django, Django REST Framework (DRF) для создания API.
\item
  База данных: PostgreSQL для хранения структурированных
  данных.
\item
  ML-модуль: (укажите используемые библиотеки или фреймворки,
  например, TensorFlow, Scikit-learn).
\item
  Дополнительные инструменты: Docker для контейнеризации, Git
  для контроля версий, CI/CD-инструменты для автоматизации
  развертывания.
\end{itemize}

\hypertarget{ux441ux445ux435ux43cux430-ux431ux430ux437ux44b-ux434ux430ux43dux43dux44bux445}{%
\subsection*{Схема базы
данных}\label{ux441ux445ux435ux43cux430-ux431ux430ux437ux44b-ux434ux430ux43dux43dux44bux445}}

База данных системы включает следующие основные таблицы: 
\begin{enumerate}
    \item 
Movies: Содержит информацию о фильмах (название, описание, год
выпуска и т.д.). 
\item 
Genres: Список жанров, к которым могут
относиться фильмы. 
\item 
Directors: Информация о режиссерах. 
\item 
Tags: Теги для фильмов, которые помогают в классификации и
поиске. 
\item 
Ratings: Рейтинги, оставленные пользователями. 
\item 
Comments: Комментарии пользователей к фильмам. 
\item 
Recommendations: Таблица для хранения рекомендаций, как для
пользователей, так и для фильмов.
\end{enumerate}
\hypertarget{uml-ux434ux438ux430ux433ux440ux430ux43cux43cux44b}{%
\subsection*{UML-диаграммы}\label{uml-ux434ux438ux430ux433ux440ux430ux43cux43cux44b}}

Для наглядности архитектуры системы можно использовать UML-диаграммы: 
\begin{itemize}
    \item 
Диаграмма классов: Отображает основные классы системы и их
взаимосвязи. 
    \item 
Диаграмма последовательностей: Показывает, как
взаимодействуют компоненты системы при выполнении типичных операций,
таких как добавление фильма или получение рекомендаций.
\end{itemize}
\hypertarget{ux443ux441ux442ux430ux43dux43eux432ux43aux430-ux438-ux43dux430ux441ux442ux440ux43eux439ux43aux430}{%
\section{Установка и
настройка}\label{ux443ux441ux442ux430ux43dux43eux432ux43aux430-ux438-ux43dux430ux441ux442ux440ux43eux439ux43aux430}}

\hypertarget{ux442ux440ux435ux431ux43eux432ux430ux43dux438ux44f-ux43a-ux43eux43aux440ux443ux436ux435ux43dux438ux44e}{%
\subsection*{Требования к
окружению}\label{ux442ux440ux435ux431ux43eux432ux430ux43dux438ux44f-ux43a-ux43eux43aux440ux443ux436ux435ux43dux438ux44e}}

Для запуска проекта необходимо следующее: 
\begin{itemize}
    \item 
Docker и
Docker Compose для контейнеризации и управления зависимостями.
    \item
Python 3.x для работы с бэкендом. 
    \item
PostgreSQL в
качестве базы данных.
\end{itemize}
\hypertarget{ux443ux441ux442ux430ux43dux43eux432ux43aux430-ux43fux440ux43eux435ux43aux442ux430}{%
\subsection*{Установка
проекта}\label{ux443ux441ux442ux430ux43dux43eux432ux43aux430-ux43fux440ux43eux435ux43aux442ux430}}

\begin{enumerate}
\def\labelenumi{\arabic{enumi}.}
\item
  Клонируйте репозиторий проекта:

\begin{Shaded}
\begin{Highlighting}[]
\FunctionTok{git}\NormalTok{ clone https://github.com/SempaiTakoo/cinema\_pon\_backend.git}
\BuiltInTok{cd}\NormalTok{ cinema\_pon\_backend}
\end{Highlighting}
\end{Shaded}
\item
  Создайте и активируйте виртуальное окружение:

\begin{Shaded}
\begin{Highlighting}[]
\ExtensionTok{python} \AttributeTok{{-}m}\NormalTok{ venv venv}
\BuiltInTok{source}\NormalTok{ venv/bin/activate  }\CommentTok{\# Для Linux/MacOS}
\ExtensionTok{venv\textbackslash{}Scripts\textbackslash{}activate}  \CommentTok{\# Для Windows}
\end{Highlighting}
\end{Shaded}
\item
  Установите зависимости:

\begin{Shaded}
\begin{Highlighting}[]
\ExtensionTok{pip}\NormalTok{ install }\AttributeTok{{-}r}\NormalTok{ requirements.txt}
\end{Highlighting}
\end{Shaded}
\item
  Запустите проект с помощью Docker:

\begin{Shaded}
\begin{Highlighting}[]
\ExtensionTok{docker{-}compose}\NormalTok{ up }\AttributeTok{{-}{-}build}
\end{Highlighting}
\end{Shaded}
\end{enumerate}

\hypertarget{ux43dux430ux441ux442ux440ux43eux439ux43aux430-ux43eux43aux440ux443ux436ux435ux43dux438ux44f}{%
\subsection*{Настройка
окружения}\label{ux43dux430ux441ux442ux440ux43eux439ux43aux430-ux43eux43aux440ux443ux436ux435ux43dux438ux44f}}

Для корректной работы системы необходимо настроить переменные окружения,
такие как: 
\begin{itemize}
    \item
DATABASE\_URL: URL для подключения к базе данных.
    \item
SECRET\_KEY: Секретный ключ для Django. 
    \item
ML\_MODEL\_PATH: Путь к модели машинного обучения.
\end{itemize}

\hypertarget{api-ux434ux43eux43aux443ux43cux435ux43dux442ux430ux446ux438ux44f}{%
\section{API
документация}\label{api-ux434ux43eux43aux443ux43cux435ux43dux442ux430ux446ux438ux44f}}

\hypertarget{ux43eux441ux43dux43eux432ux43dux44bux435-endpoints}{%
\subsection*{Основные
endpoints}\label{ux43eux441ux43dux43eux432ux43dux44bux435-endpoints}}

Система предоставляет следующие основные API endpoints: \\
Фильмы: 
\begin{itemize}
    \item GET\ /api/v1/movies/ --- получение списка.
    \item POST\ /api/v1/movies/ --- создание нового фильма.
    \item GET\ /api/v1/movies/\{id\}/ --- получение информации о
конкретном фильме.
    \item PUT\ /api/v1/movies/\{id\}/ --- обновление
информации о фильме.
    \item DELETE\ /api/v1/movies/\{id\}/ ---
удаление фильма.
\end{itemize}  

Аналогичные endpoints доступны для управления: 
\begin{itemize}
    \item 
Пользователями
(/api/v1/users/). 
    \item 
Комментариями
(/api/v1/comments/). 
    \item 
Жанрами
(/api/v1/genres/). 
    \item 
Тегами (/api/v1/tags/).
    \item 
Рейтингами (/api/v1/ratings/). 
    \item 
Режиссерами (/api/v1/directors/). 
    \item 
Рекомендациями (/api/v1/movie\_recommendations/, /api/v1/user\_recommendations/).
\end{itemize}

\hypertarget{ux43fux440ux438ux43cux435ux440ux44b-ux437ux430ux43fux440ux43eux441ux43eux432}{%
\subsection*{Примеры
запросов}\label{ux43fux440ux438ux43cux435ux440ux44b-ux437ux430ux43fux440ux43eux441ux43eux432}}

Пример запроса для получения списка фильмов:

\begin{Shaded}
\begin{Highlighting}[]
\ExtensionTok{curl} \AttributeTok{{-}X}\NormalTok{ GET http://localhost:8000/api/v1/movies/}
\end{Highlighting}
\end{Shaded}

Пример запроса для создания нового фильма:

\begin{Shaded}
\begin{Highlighting}[]
\ExtensionTok{curl} \AttributeTok{{-}X}\NormalTok{ POST http://localhost:8000/api/v1/movies/ }\AttributeTok{{-}H} \StringTok{"Content{-}Type: application/json"} \AttributeTok{{-}d} \StringTok{\textquotesingle{}\{"title": "New Movie", "genre": "Action", "director": "John Doe"\}\textquotesingle{}}
\end{Highlighting}
\end{Shaded}

\hypertarget{ux43eux440ux433ux430ux43dux438ux437ux430ux446ux438ux44f-ux43fux43eux434ux434ux435ux440ux436ux43aux438-ux438-ux441ux43eux43fux440ux43eux432ux43eux436ux434ux435ux43dux438ux44f-ux440ux430ux437ux440ux430ux431ux43eux442ux43aux438}{%
\section{Организация поддержки и сопровождения
разработки}\label{ux43eux440ux433ux430ux43dux438ux437ux430ux446ux438ux44f-ux43fux43eux434ux434ux435ux440ux436ux43aux438-ux438-ux441ux43eux43fux440ux43eux432ux43eux436ux434ux435ux43dux438ux44f-ux440ux430ux437ux440ux430ux431ux43eux442ux43aux438}}

\hypertarget{ux43fux440ux43eux446ux435ux441ux441-ux441ux43eux43fux440ux43eux432ux43eux436ux434ux435ux43dux438ux44f}{%
\subsection*{Процесс
сопровождения}\label{ux43fux440ux43eux446ux435ux441ux441-ux441ux43eux43fux440ux43eux432ux43eux436ux434ux435ux43dux438ux44f}}

Сопровождение системы включает следующие этапы: 
\begin{enumerate}
    \item 
Мониторинг:
Использование инструментов, таких как Prometheus и Grafana, для
отслеживания состояния системы. 
\item 
Обновления: Регулярное
обновление зависимостей и миграции базы данных. 
\item 
Обратная
связь: Сбор и анализ отзывов от пользователей для улучшения системы.
\end{enumerate}
\hypertarget{ux447ux430ux441ux442ux43e-ux437ux430ux434ux430ux432ux430ux435ux43cux44bux435-ux432ux43eux43fux440ux43eux441ux44b-faq}{%
\subsection*{Часто задаваемые вопросы
(FAQ)}\label{ux447ux430ux441ux442ux43e-ux437ux430ux434ux430ux432ux430ux435ux43cux44bux435-ux432ux43eux43fux440ux43eux441ux44b-faq}}

\begin{itemize}
\tightlist
\item
  Вопрос: Как добавить новый фильм? Ответ: Используйте
  endpoint POST\ /api/v1/movies/ с необходимыми данными.
\item
  Вопрос: Как получить рекомендации для пользователя?
  Ответ: Используйте endpoint
  GET\ /api/v1/user\_recommendations/\{user\_id\}/.
\end{itemize}

\hypertarget{ux440ux443ux43aux43eux432ux43eux434ux441ux442ux432ux43e-ux43fux43eux43bux44cux437ux43eux432ux430ux442ux435ux43bux435ux439-ux440ux430ux437ux440ux430ux431ux43eux442ux430ux43dux43dux43eux433ux43e-ux43fux440ux43eux434ux443ux43aux442ux430}{%
\section{Руководство пользователей разработанного
продукта}\label{ux440ux443ux43aux43eux432ux43eux434ux441ux442ux432ux43e-ux43fux43eux43bux44cux437ux43eux432ux430ux442ux435ux43bux435ux439-ux440ux430ux437ux440ux430ux431ux43eux442ux430ux43dux43dux43eux433ux43e-ux43fux440ux43eux434ux443ux43aux442ux430}}

\hypertarget{ux442ux438ux43fux44b-ux43fux43eux43bux44cux437ux43eux432ux430ux442ux435ux43bux435ux439-ux438-ux438ux445-ux43fux43eux43bux43dux43eux43cux43eux447ux438ux44f}{%
\subsection*{Типы пользователей и их
полномочия}\label{ux442ux438ux43fux44b-ux43fux43eux43bux44cux437ux43eux432ux430ux442ux435ux43bux435ux439-ux438-ux438ux445-ux43fux43eux43bux43dux43eux43cux43eux447ux438ux44f}}

\begin{itemize}
\tightlist
\item
  Администраторы:

  \begin{itemize}
  \tightlist
  \item
    Полный доступ ко всем данным.
  \item
    Возможность добавлять, редактировать и удалять данные.
  \end{itemize}
\item
  Пользователи:

  \begin{itemize}
  \tightlist
  \item
    Доступ к просмотру информации о фильмах и рекомендаций.
  \item
    Возможность оставлять комментарии и оценки.
  \end{itemize}
\end{itemize}

\hypertarget{customer-journey-map-cjm}{%
\subsection*{Customer Journey Map
(CJM)}\label{customer-journey-map-cjm}}

\hypertarget{ux430ux434ux43cux438ux43dux438ux441ux442ux440ux430ux442ux43eux440ux44b}{%
\paragraph{Администраторы}\label{ux430ux434ux43cux438ux43dux438ux441ux442ux440ux430ux442ux43eux440ux44b}}

\begin{enumerate}
\def\labelenumi{\arabic{enumi}.}
\tightlist
\item
  Вход в систему: Аутентификация через API.
\item
  Управление данными: Добавление, редактирование и удаление
  данных.
\item
  Мониторинг: Просмотр логов и мониторинг состояния системы.
\end{enumerate}

\hypertarget{ux43fux43eux43bux44cux437ux43eux432ux430ux442ux435ux43bux438}{%
\paragraph{Пользователи}\label{ux43fux43eux43bux44cux437ux43eux432ux430ux442ux435ux43bux438}}

\begin{enumerate}
\def\labelenumi{\arabic{enumi}.}
\tightlist
\item
  Регистрация/авторизация: Регистрация нового пользователя или
  вход в систему.
\item
  Просмотр фильмов: Поиск и просмотр информации о фильмах.
\item
  Получение рекомендаций: Просмотр персональных рекомендаций.
\item
  Взаимодействие: Оставление комментариев и оценок.
\end{enumerate}

\hypertarget{ux442ux435ux441ux442ux438ux440ux43eux432ux430ux43dux438ux435}{%
\section{Тестирование}\label{ux442ux435ux441ux442ux438ux440ux43eux432ux430ux43dux438ux435}}

\hypertarget{ux44eux43dux438ux442-ux442ux435ux441ux442ux44b}{%
\subsection*{Юнит-тесты}\label{ux44eux43dux438ux442-ux442ux435ux441ux442ux44b}}

Юнит-тесты покрывают основные функции системы, такие как добавление
фильмов, управление пользователями и генерация рекомендаций. Для запуска
тестов используйте команду:

\begin{Shaded}
\begin{Highlighting}[]
\ExtensionTok{python}\NormalTok{ manage.py test}
\end{Highlighting}
\end{Shaded}

\hypertarget{ux438ux43dux442ux435ux433ux440ux430ux446ux438ux43eux43dux43dux44bux435-ux442ux435ux441ux442ux44b}{%
\subsection*{Интеграционные
тесты}\label{ux438ux43dux442ux435ux433ux440ux430ux446ux438ux43eux43dux43dux44bux435-ux442ux435ux441ux442ux44b}}

Интеграционные тесты проверяют взаимодействие между различными
компонентами системы, такими как API и база данных.

\hypertarget{ux440ux430ux437ux432ux435ux440ux442ux44bux432ux430ux43dux438ux435}{%
\section{Развертывание}\label{ux440ux430ux437ux432ux435ux440ux442ux44bux432ux430ux43dux438ux435}}

\hypertarget{ux440ux430ux437ux432ux435ux440ux442ux44bux432ux430ux43dux438ux435-ux43dux430-ux441ux435ux440ux432ux435ux440ux435}{%
\subsection*{Развертывание на
сервере}\label{ux440ux430ux437ux432ux435ux440ux442ux44bux432ux430ux43dux438ux435-ux43dux430-ux441ux435ux440ux432ux435ux440ux435}}

Для развертывания системы на production-сервере используйте Docker и
CI/CD-инструменты, такие как GitHub Actions или GitLab CI.

\hypertarget{cicd}{%
\subsection*{CI/CD}\label{cicd}}

Процесс Continuous Integration/Continuous Deployment включает
автоматическое тестирование и развертывание при каждом изменении кода.

\hypertarget{ux43fux43eux434ux434ux435ux440ux436ux43aux430-ux438-ux441ux43eux43fux440ux43eux432ux43eux436ux434ux435ux43dux438ux435}{%
\section{Поддержка и
сопровождение}\label{ux43fux43eux434ux434ux435ux440ux436ux43aux430-ux438-ux441ux43eux43fux440ux43eux432ux43eux436ux434ux435ux43dux438ux435}}

\hypertarget{ux43bux43eux433ux438ux440ux43eux432ux430ux43dux438ux435}{%
\subsection*{Логирование}\label{ux43bux43eux433ux438ux440ux43eux432ux430ux43dux438ux435}}

Система использует централизованное логирование для отслеживания ошибок
и анализа производительности.

\hypertarget{ux43eux431ux440ux430ux431ux43eux442ux43aux430-ux43eux448ux438ux431ux43eux43a}{%
\subsection*{Обработка
ошибок}\label{ux43eux431ux440ux430ux431ux43eux442ux43aux430-ux43eux448ux438ux431ux43eux43a}}

Ошибки обрабатываются с использованием механизмов Django и DRF, с
возвращением соответствующих HTTP-статусов и сообщений об ошибках.

\endinput